% Paper 1, Section 1. Requires % ---------------------------------------------------------------------------
% Headline numbers used in the text. PROVISIONAL (11 Oct 2026): from the
% 657-run file after Night 4. Replace with the audited values after Nights 5-6
% (catalog v2 re-runs). Change them here only, never inline.
% ---------------------------------------------------------------------------
\newcommand{\nProblems}{73}
\newcommand{\nApps}{three}
\newcommand{\accOursEight}{49\%}      % rule v2, qwen3:8b
\newcommand{\accOursFourteen}{52\%}   % rule v2, qwen3:14b
\newcommand{\accReactEight}{12\%}     % plain ReAct agent, qwen3:8b
\newcommand{\accReactFourteen}{18\%}  % plain ReAct agent, qwen3:14b
\newcommand{\accChooseEight}{42\%}    % LLM chooses the checks, qwen3:8b
\newcommand{\tokOursEight}{5.1k}      % mean LLM tokens (in+out) per incident
\newcommand{\tokReactEight}{44k}
\newcommand{\tokOursFourteen}{4.0k}
\newcommand{\tokReactFourteen}{62k}
\newcommand{\tokChooseEight}{15k}
\newcommand{\accOursVOneFourteen}{26\%} % rule v1, qwen3:14b
\newcommand{\rqTwoWins}{three of four}  % RQ2: impact levels where ours has the lowest cost
 in the preamble (provisional values),
% % Figure 1: the incident-response lifecycle and the decisions an agent must take.
% Needs \usepackage{tikz} and \usetikzlibrary{arrows.meta,positioning,fit} in the preamble.
\begin{figure}[t]
\centering
\resizebox{\columnwidth}{!}{%
\begin{tikzpicture}[
  node distance=4mm and 3mm,
  stage/.style={draw, rounded corners=2pt, minimum height=8mm, minimum width=15mm,
                align=center, font=\footnotesize, fill=gray!8},
  dec/.style={draw, rounded corners=2pt, align=center, font=\scriptsize, fill=blue!6},
  arr/.style={-{Stealth[length=2mm]}, thick}]
  \node[stage] (det) {Detect\\\scriptsize alert / SLO breach};
  \node[stage, right=of det] (tri) {Triage\\\scriptsize is it real?};
  \node[stage, right=of tri] (dia) {Diagnose\\\scriptsize which service, why?};
  \node[stage, right=of dia] (mit) {Mitigate\\\scriptsize rollback, restart};
  \node[stage, right=of mit] (rec) {Recover \&\\\scriptsize post-mortem};
  \foreach \a/\b in {det/tri, tri/dia, dia/mit, mit/rec} \draw[arr] (\a) -- (\b);
  % the decisions inside the loop, which this paper studies
  \node[dec, below=6mm of dia, xshift=-9mm] (chk) {\textbf{check}\\next probe?};
  \node[dec, right=2mm of chk] (stp) {\textbf{stop}\\enough evidence?};
  \node[dec, right=2mm of stp] (act) {\textbf{act or}\\\textbf{escalate}?};
  \node[draw, dashed, rounded corners, fit=(chk)(stp)(act), inner sep=2mm,
        label={[font=\scriptsize]below:who decides: the LLM, or a fixed decision rule?}] (box) {};
  \draw[arr, dashed] (dia.south) -- (box.north -| dia.south);
  \draw[arr, dashed] (act.east) -| (mit.south);
\end{tikzpicture}}
\caption{The incident-response lifecycle for an operational outage. Most time is
spent in triage and diagnosis, where the responder repeatedly decides which check
to run next, when the evidence suffices, and whether to apply a reversible fix or
escalate to a human. This paper studies \emph{who} should make these decisions
when an LLM is in the loop.}
\label{fig:ir}
\end{figure}
 where Fig. 1 should float, and refs.bib.
% Length: about one page in a two-column IEEE/ACM layout, plus Fig. 1.

\section{Introduction}
\label{sec:intro}

% --- what incident response is ------------------------------------------------
Modern online services run as many interdependent microservices on shared cloud
infrastructure, and they fail in ways no single component explains. A bad
deployment, a misconfigured port, an exhausted connection pool or a pod stuck on an
unschedulable node can each surface only as rising latency or errors somewhere
downstream. \emph{Incident response} (IR) is the process by which operators
restore such a service: an alert or a breached service-level objective (SLO)
\emph{detects} the problem, responders \emph{triage} whether it is real,
\emph{diagnose} which component failed and why, \emph{mitigate} it (for example
by rolling back a release or restarting a workload), and finally recover and
write a post-mortem~\cite{beyer2016sre,nist80061r3}. Fig.~\ref{fig:ir} shows this
lifecycle. Its most expensive part is the diagnostic loop: under time pressure,
the responder repeatedly decides which check to run next, when the evidence is
sufficient, and whether to act or to call in someone more senior.

% --- recent developments -------------------------------------------------------
Large language models (LLMs) have rapidly entered this loop. Early work used LLMs
to recommend root causes and mitigations from incident text~\cite{ahmed2023recommending,chen2024rcacopilot};
tool-using agents now query logs, metrics and the cluster
directly~\cite{wang2024rcagent}, usually following the ReAct pattern of interleaved
reasoning and actions~\cite{yao2023react}. Benchmarks such as
AIOpsLab~\cite{chen2025aiopslab}, ITBench~\cite{jha2025itbench} and
OpenRCA~\cite{xu2025openrca} now inject real faults into running systems and grade
agents automatically, and several vendors market autonomous ``AI SRE'' agents.

% --- limitations / motivation ------------------------------------------------------
These benchmarks also expose the limits of the prevailing design, in which the LLM
controls the entire investigation. Success rates remain low: on ITBench, agents built on state-of-the-art models
resolve only 11.4\% of SRE scenarios~\cite{jha2025itbench}.
Because the LLM chooses its own next step, it repeats checks, wanders, and stops
when its narrative feels complete rather than when the evidence is; its stated
confidence is not calibrated against outcomes; and every step consumes tokens.
None of these behaviours is auditable in advance, which is precisely what an
operator needs before granting an agent the right to change production. Industry
analysts accordingly expect many agentic AI projects to be abandoned before they
deliver value~\cite{gartner2025agentic}. The underlying decision problem is not new:
\emph{decision-theoretic troubleshooting} chose the next observation, repair or
escalation by expected cost three decades ago~\cite{heckerman1995decision,breese1996decision,langseth2003decision},
and active probing selected the most informative test in distributed
systems~\cite{rish2005adaptive}. Those methods, however, required a probabilistic
model of each system to be built by hand, which cannot keep pace with
infrastructure that changes weekly. LLM agents removed the hand-built model, and
with it the decision theory.

% --- research questions --------------------------------------------------------------
We ask whether the two can be recombined, with the LLM building the model and a
fixed rule taking the decisions. Concretely:
\textbf{RQ1}: if an LLM proposes candidate root causes, each with the outcome it
predicts for every available check, can a fixed expected-utility rule choose which
checks to run and when to stop more efficiently than an LLM making those choices,
without losing diagnostic accuracy?
\textbf{RQ2}: can the same rule decide between applying a reversible mitigation,
gathering more evidence and escalating to a human, and does its behaviour shift
sensibly as the customer impact of the incident grows?
The rationale is to separate what LLMs do well (reading heterogeneous evidence and
imagining causes) from what they do poorly (disciplined, cost-aware sequential
decisions), and to make the latter deterministic and inspectable.

% --- proposed approach -------------------------------------------------------------
Our agent keeps an explicit belief over the LLM's hypotheses plus a ``no fault''
hypothesis and per-service catch-all causes, and updates it by Bayes' rule with a
per-check error model. For every unrun check it computes the value of information,
the expected reduction in uncertainty minus a cost term, runs the best check, and
stops when no check is worth its cost. The LLM is called only to propose and, if
every hypothesis is contradicted, to revise hypotheses; it never selects checks
or decides to stop. For mitigation, the same expected-cost calculation compares
reversible actions, further checks and escalation, with every state-changing
action additionally gated by an external policy engine.

% --- evaluation / why we trust the result -------------------------------------------------
We evaluate RQ1 on all \nProblems{} Kubernetes problems of AIOpsLab~\cite{chen2025aiopslab}
across \nApps{} applications, covering detection, localization and root-cause
analysis, with two open models (Qwen3 8B and 14B~\cite{yang2025qwen3}). Six
conditions share the same model, tools, problems and step budget and differ only in
who decides: a plain ReAct agent; the LLM choosing checks from our hypotheses;
random checks with our stopping rule; our rule; and running every check.
Comparisons are paired on identical problems with exact McNemar
tests~\cite{mcnemar1947note} and Wilson intervals~\cite{wilson1927probable}. We
audited every run for tool failures and re-ran all runs affected by a measurement
defect we found, and we report a failure of our original rule on the larger model
together with its correction. We evaluate RQ2 in a simulator of paired incidents
against a utility gate without an information term, a confidence-threshold
escalator, and a diagnose-first policy, across four levels of customer impact and
a sensitivity analysis. Our rule answers correctly on \accOursEight{} and
\accOursFourteen{} of problems (8B and 14B), against \accReactEight{} and
\accReactFourteen{} for the ReAct agent, using an order of magnitude fewer tokens
(\tokOursEight{} vs.\ \tokReactEight{} per incident on 8B). Against an LLM choosing
from the same hypotheses it is statistically tied in accuracy at about a third of
the tokens. In RQ2 it has the lowest expected cost at \rqTwoWins{} impact levels.

% --- contributions --------------------------------------------------------------------------
In summary, this paper makes the following contributions:
\begin{itemize}
  \item \textbf{An architecture} that couples LLM-generated hypothesis models with
        classical decision-theoretic troubleshooting: a fixed, inspectable rule
        selects checks, stops, and chooses between reversible action and
        escalation, while the LLM only proposes and revises hypotheses.
  \item \textbf{A controlled evaluation on a public benchmark} (\nProblems{}
        AIOpsLab problems, two models, six conditions) that isolates \emph{who
        decides} from \emph{what the model knows}, with paired statistics, cost in
        tokens and checks, and an audit of measurement errors.
  \item \textbf{Evidence that the rule matches LLM-directed accuracy at a fraction
        of the cost} and far exceeds a plain ReAct agent, together with a reported
        failure mode of the original rule on a larger model and its correction.
  \item \textbf{An act-or-escalate analysis} showing that the same rule adapts its
        willingness to act to customer impact and stays cheaper than
        confidence-threshold and information-free baselines under misspecified
        assumptions.
  \item \textbf{A reproducible artefact}, to be released with the paper: the
        agent, the benchmark harness, an offline replay tool and all run traces.
\end{itemize}
